% 原书第 9.4 节。
\section{评估基准}
\label{sec:ch09-benchmarks}
本节概述用于评估多模态大语言模型在不同模态与任务上表现的部分基准数据集，着重介绍任务种类较多、覆盖模态较广或在本领域广泛使用的基准。更详细的列表见 Yin et al. (2023)。

\begin{enumerate}
\item \textbf{CMMU} 收录 12,000 道多模态题目，由人工从大学考试、测验和教材中整理，涉及艺术与设计、商业、科学、健康与医学、人文与社会科学、技术与工程六个基础学科领域（\refcite{ch09-ref-071}{Zhang et al., 2024b}）。对应的 MMMU 则覆盖 30 个科目。CMMU 含有 39 种图像，包括统计图、示意图、地图、表格、乐谱和化学结构等，用于考察不同方面的多模态理解能力。

\item \textbf{MMCBench} 提供评估大型多模态模型（LMM）的详细框架，主要考察其在常见输入扰动下的稳健性和自洽性（\refcite{ch09-ref-073}{Zhang et al., 2024c}）。它关注文本、图像和语音模态的相互作用，涉及文本生成图像、图像生成文本、文本生成语音和语音转文本等主要任务。

\item \textbf{MMVP} 通过视觉问答任务评估多模态大语言模型的视觉能力（\refcite{ch09-ref-051}{Tong et al., 2024}）。它含有一个包含 300 张测试图像的目录，以及一个包含问题和正确答案的 CSV 文件。

\item \textbf{TimeIT} 涉及六种与时间戳相关的视频任务，纳入 12 个来自不同领域的数据集（\refcite{ch09-ref-043}{Ren et al., 2023}）。它侧重时间敏感的长视频理解任务，包括密集视频描述、视频时刻定位、视频摘要、精彩片段检测、步骤定位和语音转录，训练数据共 124,861 个实例。

\item \textbf{ViP-Bench} 包含 303 个来自 MM-Vet、MMBench 和 Visual Genome 的图像与问题配对，旨在测试多模态模型的视觉推理能力（\refcite{ch09-ref-004}{Cai et al., 2023}）。它在图像区域层面考察视觉理解的六个方面，即识别、光学字符识别（OCR）、知识、数学、物体关系推理和语言生成。基准用 GPT-4 对模型响应给出 0 至 10 分，便于定量比较。

\item \textbf{M3DBench} 汇编了超过 320K 对三维多模态指令遵循数据，包括具有不同多模态提示的 138K 多条指令（\refcite{ch09-ref-004}{Cai et al., 2023}）。它利用已有数据集以及由大语言模型为各种三维任务生成的指令。覆盖从物体检测到问答的任务，每类任务有专门的指令和响应。它通过模式匹配过滤不相关的响应，以改善数据质量，供三维指令遵循评估使用。

\item \textbf{Video-Bench} 提供评估视频大语言模型的综合基准（\refcite{ch09-ref-041}{Ning et al., 2023}）。它包含十项任务，用于考察对视频内容的理解、借助已有知识回答问题以及理解和决策等能力，共约 15,033 个实例。

\item \textbf{Bingo} 对多模态理解中模型失败与成功的实例进行分类，包括 190 个失败实例和 131 个成功实例（\refcite{ch09-ref-013}{Cui et al., 2023}）。每个实例配有一两个问题，归入“干扰”（图像对图像、文本对图像）或“偏见”（区域、OCR、事实）类别。该基准旨在分析响应出现幻觉的具体原因，用文化、语言与事实背景各异的图像，详细考察 GPT-4V(ision) 中的偏见。

\item \textbf{MMHAL-BENCH} 使用专门的指标评估多模态模型的幻觉（\refcite{ch09-ref-049}{Sun et al., 2023}）。它包含 96 个图像与问题配对，涉及八类问题和十二种物体主题，用来考察大型多模态模型是否会对图像内容作出不实陈述。为避免数据泄漏，基准采用 Open Images 验证集和测试集中的图像，覆盖“配件”“动物”“交通工具”等物体大类。GPT-4 比较模型的输出、人工编写的答案与图像内容，判断是否出现幻觉。

\item \textbf{Sparkles} 利用 GPT-4 构建多模态对话数据集，模拟围绕图像和文本展开的真实对话，涵盖多种现实场景（\refcite{ch09-ref-024}{Huang et al., 2023}）。生成过程采用两轮对话模式，先由用户就图像提问，再由助手详细回答，然后引入新图像供后续讨论。数据生成时用对话示例进行上下文学习，并根据候选图像的描述选择相关图像。当时的 GPT-4 仅接受文本输入，因此用详细的文本描述代表图像。

\item \textbf{SciGraphQA} 是一个围绕学术图表构建的大规模合成多轮问答数据集（\refcite{ch09-ref-033}{Li and Tajbakhsh, 2023}）。它含有 295,000 个样本，源自 arXiv 上 2010 至 2020 年的 290,000 篇计算机科学或机器学习论文。用 PaLM-2 根据论文中的图表生成对话，并纳入标题、摘要、相关段落及上下文数据。每段对话平均有 2.23 轮问答。在 3,000 个样本的测试集上，GPT-4 给出的问答匹配质量平均为 8.7 分，满分为 10 分。

\item \textbf{LAMM} 是一个多模态指令调优数据集，用多种视觉任务中的图像与点云，构建 186,000 对语言与图像、10,000 对语言与三维数据的指令和响应（\refcite{ch09-ref-064}{Yin et al., 2024}）。数据集利用 GPT API 和自指令方法构建，包括日常对话、事实知识对话、图像与三维场景的详细描述、视觉任务对话四类配对，以增强视觉任务泛化能力。它涉及多种二维和三维视觉任务，例如描述生成、场景图识别、视觉问答、分类、检测、计数和 OCR。
\end{enumerate}
